\begin{proof}
\textbf{Theorem (Jensen).}  
Let $f:\mathbb R\to\mathbb R$ be convex, i.e.
\[
f\big(t x+(1-t)y\big)\le t\,f(x)+(1-t)f(y)\qquad\text{for all }x,y\in\mathbb R,\;t\in[0,1].
\tag{C}
\]
For any integer $n\ge1$, real numbers $x_{1},\dots,x_{n}$ and non‑negative
weights $\lambda_{1},\dots,\lambda_{n}$ with $\sum_{i=1}^{n}\lambda_{i}=1$, one has
\[
f\!\Bigl(\sum_{i=1}^{n}\lambda_{i}x_{i}\Bigr)\le
\sum_{i=1}^{n}\lambda_{i}f(x_{i}).\tag{J}
\]

\medskip
\textbf{Proof by induction on $n$.}

\textbf{1. Degenerate case.}  
If $\lambda_{k}=1$ for some index $k$ and $\lambda_{i}=0$ for $i\neq k$, then both sides of \eqref{J} equal $f(x_{k})$; hence the inequality holds with equality.

\medskip
\textbf{2. Base case $n=2$.}  
Let $\lambda_{1},\lambda_{2}\ge0$ with $\lambda_{1}+\lambda_{2}=1$. Applying \eqref{C} with $t=\lambda_{1}$ gives
\[
f(\lambda_{1}x_{1}+\lambda_{2}x_{2})
\le \lambda_{1}f(x_{1})+\lambda_{2}f(x_{2}),
\]
which is exactly \eqref{J} for two points.

\medskip
\textbf{3. Inductive hypothesis.}  
Assume that for some $k\ge2$ the inequality \eqref{J} holds for every collection of $k$ points with non‑negative weights summing to $1$, i.e.
\[
f\!\Bigl(\sum_{i=1}^{k}\lambda_{i}x_{i}\Bigr)
   \le \sum_{i=1}^{k}\lambda_{i}f(x_{i})\tag{IH}
\]
for all $\lambda_{i}\ge0$ with $\sum_{i=1}^{k}\lambda_{i}=1$.

\medskip
\textbf{4. Inductive step $k\to k+1$.}  
Let $\lambda_{1},\dots,\lambda_{k+1}\ge0$ with $\sum_{i=1}^{k+1}\lambda_{i}=1$ and set $\mu:=\lambda_{k+1}$. If $\mu=1$ we are in the degenerate situation of Step 1, so assume $\mu\neq1$.

Define new weights for the first $k$ points
\[
\alpha_{i}:=\frac{\lambda_{i}}{1-\mu},\qquad i=1,\dots,k .
\]
Because $\sum_{i=1}^{k}\lambda_{i}=1-\mu$, we have $\alpha_{i}\ge0$ and
\[
\sum_{i=1}^{k}\alpha_{i}=1 .
\]

Hence the full convex combination can be written as a two‑point mixture:
\begin{align*}
\sum_{i=1}^{k+1}\lambda_{i}x_{i}
   &=(1-\mu)\Bigl(\sum_{i=1}^{k}\alpha_{i}x_{i}\Bigr)+\mu x_{k+1}. \tag{4.1}
\end{align*}
Applying convexity \eqref{C} to the right‑hand side of \eqref{4.1} with
$t=1-\mu$ yields
\begin{align*}
f\!\Bigl(\sum_{i=1}^{k+1}\lambda_{i}x_{i}\Bigr)
   &\le (1-\mu)\,
      f\!\Bigl(\sum_{i=1}^{k}\alpha_{i}x_{i}\Bigr)+\mu f(x_{k+1}). \tag{4.2}
\end{align*}
By the inductive hypothesis \eqref{IH} applied to the inner combination
$\sum_{i=1}^{k}\alpha_{i}x_{i}$ we have
\begin{align*}
f\!\Bigl(\sum_{i=1}^{k}\alpha_{i}x_{i}\Bigr)
   &\le \sum_{i=1}^{k}\alpha_{i}f(x_{i}). \tag{4.3}
\end{align*}
Substituting \eqref{4.3} into \eqref{4.2} gives
\begin{align*}
f\!\Bigl(\sum_{i=1}^{k+1}\lambda_{i}x_{i}\Bigr)
   &\le (1-\mu)\sum_{i=1}^{k}\alpha_{i}f(x_{i})+\mu f(x_{k+1}). \tag{4.4}
\end{align*}
Using the definition of $\alpha_{i}$ and the identity $(1-\mu)\alpha_{i}=\lambda_{i}$,
the first term on the right‑hand side of \eqref{4.4} simplifies:
\begin{align*}
(1-\mu)\sum_{i=1}^{k}\alpha_{i}f(x_{i})
   &=\sum_{i=1}^{k}\lambda_{i}f(x_{i}). \tag{4.5}
\end{align*}
Consequently,
\begin{align*}
f\!\Bigl(\sum_{i=1}^{k+1}\lambda_{i}x_{i}\Bigr)
   &\le \sum_{i=1}^{k}\lambda_{i}f(x_{i})+\mu f(x_{k+1})\\
   &=\sum_{i=1}^{k+1}\lambda_{i}f(x_{i}),
\end{align*}
which is precisely Jensen’s inequality for $k+1$ points.

Thus the statement holds for $k+1$ whenever it holds for $k$. Together
with the base case $n=2$ (and the trivial case $n=1$) the inequality
\eqref{J} is true for every integer $n\ge1$.

\medskip
\textbf{5. Equality condition.}  
Equality in the final inequality can occur only when equality holds in
every step used above.

\begin{itemize}
\item In the outer application of convexity \eqref{4.2}, equality requires
either $\mu\in\{0,1\}$ (the mixture collapses to a single point) \emph{or}
\[
\sum_{i=1}^{k}\alpha_{i}x_{i}=x_{k+1}.
\]
\item In the inner inequality \eqref{4.3}, equality holds precisely when
the restriction of $f$ to the convex hull of $\{x_{1},\dots,x_{k}\}$ is
affine; for a strictly convex $f$ this forces all those points to be
equal (or all but one weight zero).
\end{itemize}
Hence, in general, equality in Jensen’s inequality holds \emph{iff}
the points with positive weights lie in a region where $f$ is affine.
In particular, if $f$ is \emph{strictly} convex, equality occurs only in
the two situations listed in Step 1:

\begin{enumerate}
\item one weight equals $1$ (the degenerate case), or
\item all points $x_{1},\dots,x_{n}$ are equal.
\end{enumerate}

\medskip
\textbf{6. Remarks on zero weights.}  
If some $\lambda_{i}=0$ the corresponding term disappears from both sides
of \eqref{J}; the argument above is unchanged, so the inequality remains
valid.
\end{proof}